\section{Кодируемый объект и постановка задачи}
\label{sec:setup}

\subsection{Разные значения длины описания}

Пусть наблюдения $Y_1,Y_2,\dots$ независимы и имеют распределение $Q$,
а $P_n$ --- предиктор после $n$ обучающих примеров.
Все дивергенции и натуральные логарифмы измеряются в натах;
$\log_2$ и длины бинарных сообщений --- в битах. Положим
\[
 K_n=\KL(Q\Vert P_n),\qquad S(N)=\sum_{n=0}^{N-1}K_n.
\]
Для случайной истории $K_n$ также случайна.

Если байесовский предиктор обновляется после каждого наблюдения, то цепное
правило даёт точное равенство
\begin{equation}
 \E[-\ln P(Y_{1:N})]=NH(Q)+\sum_{n=0}^{N-1}\E K_n.
 \label{eq:prequential}
\end{equation}
Здесь $H(Q)$ --- энтропия одного дискретного наблюдения в натах, а
$P(Y_{1:N})$ --- байесовская смесь. Избыточность правой части не следует
отождествлять с длиной отдельного сообщения о весах.

В регулярных моделях при стандартных условиях ведущий коэффициент байесовской
избыточности равен половине размерности. Для параметрического семейства
$p_\theta$ с априорной мерой $\Pi$ определим эмпирическую потерю и байесовскую
свободную энергию:
\[
 L_N(\theta)=-\frac1N\sum_{i=1}^{N}\ln p_\theta(Y_i),\qquad
 F_N=-\ln\int e^{-NL_N(\theta)}\,\Pi(d\theta).
\]
Пусть $Q$ реализуется параметром $\theta_*$ этого семейства. В сингулярной
модели при условиях аналитичности и интегрируемости сингулярной теории обучения
в разложении $F_N$ появляются действительный лог-канонический порог (RLCT)
$\lambda$ и кратность $m$ \citep{watanabe2009,watanabe2013}:
\begin{equation}
 F_N=NL_N(\theta_*)
       +\lambda\ln N-(m-1)\ln\ln N+O_{\Prob}(1).
 \label{eq:slt}
\end{equation}
Здесь $O_{\Prob}(1)$ обозначает ограниченный по вероятности остаток.
Это разложение характеризует байесовское интегрирование по параметрам;
асимптотика отдельной оптимизационной траектории требует самостоятельного
анализа. Связь RLCT с практической сжимаемостью исследуется, в частности,
в \citet{urdshals2025}.

\subsection{Requential-протокол}

На каждом шаге кодер и декодер знают состояние ученика, правило обновления и
размер следующего пакета $B$. Кодер использует учителя $Q$ для выбора пакета
из общего потока кандидатов; декодеру достаточно полученного индекса.
В основе выбора лежит пуассоновское функциональное представление (PFR)
\citep{li2018}. Для распределения предложения $P$ и $Q\ll P$ положим
$R=dQ/dP$. Пусть $T_1<T_2<\dots$ --- времена прихода пуассоновского процесса
единичной интенсивности, а $X_1,X_2,\dots$ --- независимые метки с законом $P$,
не зависящие от времён. Тогда
\[
 J=\argmin_{j\ge1}\frac{T_j}{R(X_j)},\qquad X_J\sim Q,
\]
где при $R(X_j)=0$ отношение принимается равным $+\infty$.
Процесс времён и меток составляет общую случайность кодера и декодера;
для разных пакетов используются её независимые части.

В основном протоколе распределение предложения для пакета равно
$P_n^{\otimes B}$: состояние ученика заморожено до конца пакета. Распределение
выбранного пакета равно $Q^{\otimes B}$. Индекс $J\ge1$ передаётся отдельным
$\delta$-кодом Элиаса \citep{elias1975}, имеющим длину
\begin{equation}
 \ell_\delta(j)=\lfloor\log_2j\rfloor
  +2\left\lfloor\log_2(\lfloor\log_2j\rfloor+1)\right\rfloor+1.
 \label{eq:elias}
\end{equation}
В частности, $\ell_\delta(1)=1$: даже совпадение распределений не устраняет
обязательное сообщение. Для пакетной дивергенции $h=BK_n$ используется
известная условная верхняя оценка \citep[приложение A]{qiu2026}
\begin{equation}
 \E[\ell_\delta(J)\mid\text{история}]
 \le f(h/\ln2),\qquad
 f(d)=d+2\log_2(1+d)+\kappa,
 \quad 0<\kappa<5.21.
 \label{eq:universal-bound}
\end{equation}
Здесь используется константа
$\kappa=1+\beta+2\log_2(1+\beta)$ при $\beta=1+(e\ln2)^{-1}$.
Правая часть оценивает условное математическое ожидание длины сообщения.

Фиксируем заранее детерминированное расписание
$0=n_0<n_1<\dots<n_m=N$, $B_i=n_{i+1}-n_i$.
Правило остановки, $N$, априорное распределение, правило обновления и код
известны обеим сторонам;
их описание в рассматриваемую условную длину не входит. Расписание,
подобранное после наблюдения неизвестных декодеру данных, требует отдельной
передачи. В разделе~\ref{sec:batch} учитель детерминирован, поэтому все
состояния и оптимальные расписания вычисляются по публичным параметрам.

\subsection{Отличие от существующих постановок}

Работы о prequential-кодах исследуют последовательное обучение, replay и
калибровку \citep{blier2018,bornschein2022}. В \citet{finzi2026} requential-код
вводится как способ описать усвоенную структуру. В \citet{qiu2026} развито
сжатие модели через выбранные данные; в формуле (21) рассматривается разбиение
\emph{уже заданного} обучающего пакета на сообщения при сохранении исходных
обновлений. Здесь размер пакета определяет следующее обновление ученика и
потому меняет его будущие предсказания. Наша оптимизация проводится при
фиксированном конечном риске.

Анализ линейных факторизаций связан с работами об инвариантах баланса и
неявном смещении градиентного потока \citep{du2018,tarmoun2021,min2022}.
Мы рассматриваем дискретный спуск с гармонически убывающим шагом и вычисляем
асимптотику \emph{суммы} рисков. Восстановление конечной точки связано с
последовательным обучением управляемым учителем \citep{lessard2019}; цена
управляющего действия у нас задаётся длиной индекса в общем потоке.
